% TOP_PROBLEM: {"id": 8000007, "title": "Non-self adjoint spectral problems", "category_id": 16, "category": {"id": 16, "name": "physics", "display_name": "Mathematical Physics", "description": "Problems at the intersection of mathematics and physics.", "slug": "physics", "order_index": 16, "created_at": "2026-07-31T15:26:25.671Z"}, "status": "partially_solved", "statusQualification": "The source’s programme includes eigenvalue accumulation location and rate for general-data AKNS operators. Its informal assertion about pseudomodes throughout the numerical range is source context, not an added universal theorem for arbitrary operators. The source’s distance wording is corrected to reciprocal distance: for self-adjoint A the resolvent norm is the reciprocal of the distance to its spectrum, at resolvent parameters."}
% FIRST_POSTED: 2026-10-11
% ENTRY_KIND: statement_only
% SUBMISSION: {"kind": "statement_only", "category": "top_problems", "problem_id": "8000007", "proposal_type": "established", "source": "Percy Deift, Some open problems in random matrix theory and the theory of integrable systems (2007), arXiv:0712.0849, Problem 7. https://arxiv.org/abs/0712.0849", "status_evidence": "Inherited source-catalogue status partially_solved preserved. Primary source accessed 11 October 2026; compared against the supplied catalogue excerpt. Historical wording and partial-result remarks were checked for scope; no comprehensive current-literature or proof-verification audit was performed.", "status_checked_on": "2026-10-11", "duplicate_check": "Full 15,773-record source registry checked for identical normalized statements; available site registry checked for identical normalized title/statement and canonical-ID collision. Canonical ID is absent from the site registry. Related canonical records 8100006, 20002493, 20002494. No mathematical-equivalence literature audit performed.", "attribution": "Percy Deift. Statement preparation: OpenAI GPT-6 Codex assistant, generation 11 October 2026; source-matching and notation editing by the assistant, no independent human mathematical verification. Inherited rights/provenance note: Public arXiv source TeX; arXiv reuse license NEEDS_REVIEW"}
% SOURCE_STATEMENT: "Much is now known, theoretically and also numerically, about the spectrum of self-adjoint operators. This is in great contrast to the situation regarding non-self-adjoint operators, where the spectral theory is far more subtle and numerical schemes must overcome significant, inherent instabilities. The heart of the difficulties and instabilities lies in the following fact: in the self-adjoint case $A=A^*$, the resolvent $(A-\\lambda)^{-1}$ is bounded by the distance of $\\lambda$ to the spectrum of $A$, but in the non-self-adjoint case, this is no longer true, as we see already in the $2\\times2$ case, $A=\\begin{bmatrix}0 n\\\\0 0\\end{bmatrix}$, $n\\to\\infty$. In recent years, a number of authors (e.g. Trefethan, Davies,...) have initiated a systematic approach to non-self-adjoint spectral problems, notions such as pseudospectrum have come into prominence, and other authors have conducted in depth studies of particular non-self-adjoint spectrum problems which arise in practice. For example, in analyzing the semi-classical limit of the focusing NLS equation, one must analyze the spectrum of the associated AKNS operator $T(h)$, as Planck's constant $h$ goes to zero. The operator is non-self-adjoint, and as $h$ goes to zero more and more eigenvalues, corresponding to solitons, emerge in the complex plane. It is of critical importance to the analysis of the semi-classical limits for NLS to determine where in the plane, and at what rate, the eigenvalues accumulate. The difficulty in doing this, theoretically and numerically, is illustrated by the following fact: the spectrum of $T(h)$ off the real axis is a discrete set, whereas the numerical range of $T(h)$ is an open subset of the plane. Nevertheless, every point $\\lambda$ lying in the numerical range of $T(h)$ is an eigenvalue of $T(h)$ to all orders in $h$, $\\|(T(h)-\\lambda)u\\|=O(|h|^k)$ for any $k\\geq 1$, for some $u=u(h)$, $\\|u\\|=1$. In the language of Kruskal, the computation of the spectrum of $T(h)$ is a problem ``beyond all orders.'' Much has been done (Kamvissis-Miller-McLaughlin, Tovbis-Venakides-Zhou,...) in analyzing the semi-classical limit for NLS in special situations where the spectral problem can be solved explicitly. The spectral problem with general data, both for NLS and also other related non-self-adjoint problems, however, is far from understood, and poses a great challenge whose resolution is still only in the initial stages."
% Preparation and final self-audit: CONTRIBUTING.md prompts C and D; source/audit artifacts retained in the local physics_statements_2026_10_11 batch.
% Imported from: physics_statements.tex; UnsolvedMath ID 8000007
% !TeX program = lualatex
\documentclass[11pt]{article}
\usepackage[margin=25mm]{geometry}
\usepackage{fontspec}
\setmainfont{Latin Modern Roman}
\usepackage{amsmath,amssymb,mathtools}
\usepackage{unicode-math}
\setmathfont{Latin Modern Math}
\usepackage{xurl,hyperref}
\hypersetup{colorlinks=true,urlcolor=blue}
\setcounter{secnumdepth}{0}
\setlength{\parindent}{0pt}
\setlength{\parskip}{5pt}
\setlength{\emergencystretch}{3em}
\newcommand{\sref}[2]{\hyperlink{source:#1:#2}{[S#2]}}
\newcommand{\cataloguescope}[1]{\paragraph{Recorded target.}#1}
\newcommand{\R}{\mathbb R}
\newcommand{\C}{\mathbb C}
\newcommand{\N}{\mathbb N}
\newcommand{\Z}{\mathbb Z}
\newcommand{\Q}{\mathbb Q}
\newcommand{\F}{\mathbb F}
\newcommand{\eps}{\varepsilon}
\begin{document}
% UnsolvedMath ID 8000007; source catalogue code AMR-079-0007
\setcounter{section}{8000006}
\section{Non-self adjoint spectral problems}\label{8000007}
{\small UnsolvedMath ID 8000007 · Mathematical Physics}
\subsection{Definitions and mathematical statement}

\paragraph{Shared notation.}$\mathbb N=\{1,2,\ldots\}$; the standard number fields have their usual meanings. All additional conventions are stated locally.


\paragraph{Definitions and notation.}For a closed operator $A$, the spectrum $\sigma(A)$ consists of parameters where $A-\lambda$ has no bounded inverse. Its $\varepsilon$-pseudospectrum also includes parameters where $\|(A-\lambda)^{-1}\|>\varepsilon^{-1}$. Numerical range consists of $\langle Au,u\rangle$ for unit vectors in the domain. A semiclassical pseudomode is a unit vector with residual tending rapidly to zero as $h\to0$; it need not correspond to an actual eigenvalue.

\paragraph{Catalogue wording retained for traceability.}
Much is now known, theoretically and also numerically, about the spectrum of self-adjoint operators. This is in great contrast to the situation regarding non-self-adjoint operators, where the spectral theory is far more subtle and numerical schemes must overcome significant, inherent instabilities. The heart of the difficulties and instabilities lies in the following fact: in the self-adjoint case $A=A^*$, the resolvent satisfies $\|(A-\lambda)^{-1}\|=1/\operatorname{dist}(\lambda,\sigma(A))$ for $\lambda\notin\sigma(A)$, but in the non-self-adjoint case, this is no longer true, as we see already in the $2\times2$ case, $A=\begin{bmatrix}0&n\\0&0\end{bmatrix}$, $n\to\infty$. In recent years, a number of authors (e.g. Trefethan, Davies,...) have initiated a systematic approach to non-self-adjoint spectral problems, notions such as pseudospectrum have come into prominence, and other authors have conducted in depth studies of particular non-self-adjoint spectrum problems which arise in practice. For example, in analyzing the semi-classical limit of the focusing NLS equation, one must analyze the spectrum of the associated AKNS operator $T(h)$, as Planck's constant $h$ goes to zero. The operator is non-self-adjoint, and as $h$ goes to zero more and more eigenvalues, corresponding to solitons, emerge in the complex plane. It is of critical importance to the analysis of the semi-classical limits for NLS to determine where in the plane, and at what rate, the eigenvalues accumulate. The difficulty in doing this, theoretically and numerically, is illustrated by the following fact: the spectrum of $T(h)$ off the real axis is a discrete set, whereas the numerical range of $T(h)$ is an open subset of the plane. Nevertheless, every point $\lambda$ lying in the numerical range of $T(h)$ is an eigenvalue of $T(h)$ to all orders in $h$, $\|(T(h)-\lambda)u\|=O(|h|^k)$ for any $k\geq 1$, for some $u=u(h)$, $\|u\|=1$. In the language of Kruskal, the computation of the spectrum of $T(h)$ is a problem ``beyond all orders.'' Much has been done (Kamvissis-Miller-McLaughlin, Tovbis-Venakides-Zhou,...) in analyzing the semi-classical limit for NLS in special situations where the spectral problem can be solved explicitly. The spectral problem with general data, both for NLS and also other related non-self-adjoint problems, however, is far from understood, and poses a great challenge whose resolution is still only in the initial stages.
\paragraph{Scope and source qualification.}The source’s programme includes eigenvalue accumulation location and rate for general-data AKNS operators. Its informal assertion about pseudomodes throughout the numerical range is source context, not an added universal theorem for arbitrary operators. The source’s distance wording is corrected to reciprocal distance: for self-adjoint A the resolvent norm is the reciprocal of the distance to its spectrum, at resolvent parameters.
\subsection{Short English statement}
Understand and compute non-self-adjoint spectra, especially the semiclassical scattering spectrum of focusing nonlinear Schrödinger equations.
\subsection{Sources}
\begin{itemize}
\item[\hypertarget{source:8000007:1}{S1}] ULAM.AI and the UnsolvedMath Contributors, UnsolvedMath, record 8000007 (AMR-079-0007). Catalogue curation under CC BY 4.0; definitions and English exposition added locally.
\url{https://huggingface.co/datasets/ulamai/UnsolvedMath}.
\item[\hypertarget{source:8000007:2}{S2}] Percy Deift, Some open problems in random matrix theory and the theory of integrable systems (2007), arXiv:0712.0849, Problem 7.
\url{https://arxiv.org/abs/0712.0849}.
\end{itemize}
\label{end:8000007}
\end{document}
